\textbf{第一不完全性定理}指出，对于任何一致、可公理化并且扩张~$\Th{Q}$ 的理论~$\Th{T}$，
都存在一个语句~$!G_\Th{T}$，使得 $\Th{T} \Proves/
!G_\Th{T}$。
$!G_\Th{T}$ 的构造使得它以一种迂回的方式表达"$\Th{T}$ 不能证明~$!G_\Th{T}$"。
既然 $\Th{T}$ 事实上不能证明它，那么它所说的就是真的。
如果 $\Sat{N}{\Th{T}}$，那么 $\Th{T}$ 不会证明任何假命题，因此 $\Th{T} \Proves/ \lnot !G_\Th{T}$。
这样的语句在~$\Th{T}$ 中是\textbf{独立的}或称\textbf{不可判定的}。
G\"odel 最初的证明在假定 $\Th{T}$ 是\textbf{$\omega$-一致的}的前提下，确立了 $!G_\Th{T}$ 的独立性。
Rosser（罗瑟）改进了这一结果，他找到了另一个语句~$!R_\Th{T}$，只要 $\Th{T}$ 是简单一致的，该语句在~$\Th{T}$ 中就既不可证也不可否证。

$!G_\Th{T}$ 的构造是能行的：给定~$\Th{T}$ 的一个公理化，我们原则上可以写出~$!G_\Th{T}$。
$!G_\Th{T}$ 陈述其自身不可证性的那种"迂回方式"，是一个更一般结果——\textbf{不动点引理}——的特例。
该引理断言，对于 $\Lang{L_A}$ 中的任意公式~$!B(y)$，都存在一个语句~$!A$，使得 $\Th{Q} \Proves !A \liff
!B(\gn{!A})$。
（这里，$\gn{!A}$ 是~$!A$ 的哥德尔数的标准数码，即 $\num{\Gn{!A}}$。）
为了得到 $!G_\Th{T}$，我们取公式~$\lnot\OProv[\Th{T}](y)$ 作为~$!B(y)$。
我们此前工作的最终成果便是 $\OProv[\Th{T}]$：我们知道 $\Prf[\Th{T}](n, m)$ 是原始递归的——
当 $n$ 为从 $\Th{T}$ 出发、以哥德尔数为 $m$ 的语句为结论的某条推导的哥德尔数时，该关系成立。
我们也知道 $\Th{Q}$ 能表示所有原始递归关系，因此存在某个公式 $\OPrf[\Th{T}](x,y)$ 在~$\Th{Q}$ 中表示~$\Prf[\Th{T}]$。
$\Th{T}$ 的\textbf{可证性谓词}是 $\Prov[\Th{T}](y)$，即 $\lexists[x][\Prf[\Th{T}](x, y)]$，它表达了在~$\Th{T}$ 中可证这一性质。
（但它并不表示该性质：如果 $\Th{T} \Proves !A$，那么 $\Th{Q} \Proves \OProv[\Th{T}](\gn{!A})$；
但如果 $\Th{T} \Proves/ !A$，那么 $\Th{Q}$ 一般并不证明 $\lnot
\OProv[\Th{T}](\gn{!A})$。）

\textbf{第二不完全性定理}确立了 $\Th{T}$ 也不能证明语句~$\OCon[\Th{T}]$——该语句表达 $\Th{T}$ 是一致的，也即 $\Th{T}$~不能证明 $\lnot
\OProv[\Th{T}](\gn{\lfalse})$。
第二不完全性定理的证明需要 $\Th{T}$ 满足一些额外条件，即\textbf{可证性条件}。
$\Th{PA}$ 满足这些条件，尽管~$\Th{Q}$ 不满足。
满足可证性条件的理论也满足\textbf{L\"ob（勒布）定理}：$\Th{T} \Proves
\OProv[\Th{T}](\gn{!A}) \lif !A$ 当且仅当 $\Th{T} \Proves !A$。

不动点定理还有另一个重要推论。
我们称关系 $R(n)$ 在 $\Lang{L_A}$ 中是\textbf{可定义的}，如果存在公式~$!A_R(x)$ 使得 $\Sat{N}{!A_R(\num{n})}$ 当且仅当 $R(n)$ 成立。
例如，$\Prov[\Th{T}]$ 是可定义的，因为 $\OProv[\Th{T}]$ 定义了它。
然而，"是真语句的哥德尔数"这一性质（即 $n$ 具有该性质当且仅当存在语句 $!A$ 使得 $n = \Gn{!A}$ 且 $\Sat{N}{!A}$）是不可定义的。
这就是关于真之不可定义性的\textbf{Tarski（塔斯基）定理}。